% ======================================
% Capstone Assessment: Predicate Calculus
% ======================================

\section{Capstone Assessment: Predicate Calculus}

\begin{remark*}[Assessment scope]
This capstone checks whether translations, syntactic scope, semantic
counterexamples, and proof-rule choices fit together. Use only the predicate
logic developed in this chapter: domains, translation keys, formulas,
satisfaction, consequence, quantifier rules, and equality when it is explicitly
part of the language.
\end{remark*}

\begin{enumerate}
  \item \textbf{Choose a domain and translation key.}
  Translate a short paragraph containing people, courses, and prerequisites.
  State the domain, each constant, each predicate, and each relation before
  writing formulas.

  \item \textbf{Translate English statements.}
  Include one universal claim, one existential claim, one negative claim, and
  one mixed-quantifier claim. Explain why each restriction uses implication or
  conjunction.

  \item \textbf{Identify free and bound variables.}
  For each displayed formula, list the free variables and bound variables.
  Repair any formula that was intended to be a sentence but still has a free
  variable.

  \item \textbf{Evaluate scope ambiguity.}
  Compare \(\forall x\exists y\,R(x,y)\) with
  \(\exists y\forall x\,R(x,y)\). Give an English reading of each and explain
  which one matches the paragraph.

  \item \textbf{Test a claimed inference.}
  Decide whether the proposed conclusion follows from the translated premises.
  If it is valid, give a proof using quantifier rules. If it is invalid, give a
  small counterexample by describing a domain and symbol readings.
\end{enumerate}

\begin{remark*}[Completion criterion]
Successful solutions keep the domain fixed, use one translation key
consistently, respect variable scope, distinguish dependent witnesses from
uniform witnesses, and justify invalidity with a concrete counterexample rather
than an informal objection.
\end{remark*}
